\chapter{Genitourinary Syndrome of Menopause}

\section*{Definition}
Genitourinary syndrome of menopause (GSM) is a collection of vulvovaginal, sexual, and urinary symptoms caused by reduced estrogen activity during perimenopause and after menopause. It may include dryness, burning, irritation, painful sex, urinary urgency, recurrent urinary tract infections, and changes in the vulva or vagina.

\section*{When to See a Doctor}
Arrange medical assessment for persistent dryness, burning, vulval or vaginal pain, painful sex, new urinary symptoms, recurrent urinary infections, or bleeding after sex. Urgent assessment is needed for heavy bleeding, postmenopausal bleeding, fever with flank pain, inability to urinate, severe pelvic pain, or visible vulval or vaginal lesions.

\section*{How It Develops}
Reduced estrogen causes thinning of the vaginal and vulval epithelium, reduced lubrication, loss of elasticity, decreased blood flow, and a rise in vaginal pH. The normal protective lactobacillus-dominant environment may change. Urethral and bladder tissues can also become more sensitive, contributing to urgency, dysuria, and recurrent infection.

\section*{Causes}
\begin{itemize}
  \item Perimenopause, menopause, and surgical or treatment-induced menopause.
  \item Ovarian suppression, pelvic radiotherapy, chemotherapy, or anti-estrogen treatment.
  \item Breastfeeding or other low-estrogen states can produce similar symptoms.
  \item Irritants such as perfumed washes, douches, harsh soaps, and some topical products may worsen symptoms.
\end{itemize}

\section*{Related Symptoms}
\begin{itemize}
  \item Vaginal or vulval dryness, itching, burning, irritation, reduced lubrication, and painful intercourse.
  \item Urinary urgency, frequency, dysuria, recurrent urinary tract infections, stress incontinence, and nocturia.
  \item Labial or clitoral changes, vaginal narrowing, reduced sexual comfort, and pelvic organ prolapse symptoms.
\end{itemize}

\section*{Medical Evaluation}
\begin{itemize}
  \item Assessment includes symptom pattern, menstrual and menopause history, sexual and urinary symptoms, medications, cancer treatment, and personal goals.
  \item Examination is guided by symptoms and may assess the vulva, vagina, pelvic floor, urethra, prolapse, infection, skin disease, or other lesions.
  \item Urinalysis and culture are used for urinary symptoms; persistent bleeding, unexplained lesions, or atypical findings require appropriate gynecologic evaluation.
\end{itemize}

\section*{Treatment Options}
Treatment is individualised. Vaginal moisturisers and lubricants may help mild symptoms. Low-dose vaginal estrogen is effective for many GSM symptoms and has minimal systemic absorption; alternatives include vaginal prasterone or ospemifene when appropriate. Pelvic-floor therapy, bladder training, and treatment of confirmed infection may be added. Breast-cancer history and other contraindications require shared decision-making with the relevant clinician.

\section*{Self-Care and Home Management}
Use an unscented emollient externally and a regular vaginal moisturiser if helpful. Water- or silicone-based lubricants can reduce friction during sex. Avoid douching, perfumed products, and irritant soaps. Seek assessment rather than repeatedly self-treating presumed thrush or urinary infection.

\section*{References}
\begin{thebibliography}{99}
\bibitem{gsm:ref1} The North American Menopause Society. The 2022 hormone therapy position statement. \textit{Menopause}. 2022;29:767--794.
\bibitem{gsm:ref2} National Institute for Health and Care Excellence. Menopause: diagnosis and management. NICE guideline NG23; updated 2024.
\bibitem{gsm:ref3} American College of Obstetricians and Gynecologists. Treatment of urogenital symptoms in individuals with a history of estrogen-dependent breast cancer. Clinical Consensus; 2021.
\end{thebibliography}
